\section{咸丰同治的账本叙事：战争、省份与重建}
1851年至1874年的事件需要在省级军队、河流与城市、宗教与族群分类、外交文书和战后企业之间展开。本节以账本的主证和保留为底线：战报不等于伤亡普查，条约不等于即时实施，地方社会也不能由团练、教民或回民等行政类别一笔概括。
\subsection{1850—1859：事件、文书与地方后果}
这一组节点按账本的日期顺序排列。相邻事件并不必然构成单一因果，却共同显示信息、资源与地方行动如何在同一时段交错。
\hypertarget{ledger:EQ-1851-JINTIAN-UPRISING}{}\paragraph{咸丰元年（1851年）}洪秀全集团在金田起事，太平天国战争开始。核心取证为《清文宗实录》咸丰元年相关条；军机处档、战事方略及地方奏报。编年记录适于钉住事件在朝廷文书中的时间位置，却需以地方、对方或物质材料检验其范围和后果。账本所列条件为：宗教组织、地方社会冲突和清末财政治理危机交织。其后果被限定为：参与者动机与人数不能由官府或太平军单方叙事概括。取证保留：以咸丰元年（1851年）的同期文书为定位起点；官修记录、奏报或外交文本服务特定机构立场，具体日期、规模、地方执行与对手视角须以独立材料复核。
\hypertarget{ledger:EQ-1851-YONGAN-KINGSHIPS}{}\paragraph{咸丰元年（1851年）}太平军在永安建立王号和军事组织。核心取证为《清文宗实录》咸丰元年相关条；宫中朱批奏折、内阁大库题本与《清史稿》相关志传。这类上谕、题本或部议首先证明机构处理了何种命令、呈报与责任，而不自动证明地方收到、执行或接受。账本所列条件为：起事扩展需要更稳定的指挥与合法性语言。其后果被限定为：太平天国制度在流动战时环境中不断变化。取证保留：以咸丰元年（1851年）的同期文书为定位起点；官修记录、奏报或外交文本服务特定机构立场，具体日期、规模、地方执行与对手视角须以独立材料复核。
\hypertarget{ledger:EQ-1852-TAIPING-NORTHWARD-MARCH}{}\paragraph{咸丰二年（1852年）}太平军北上经湖南，清军防御与地方团练开始加速组织。核心取证为《清文宗实录》咸丰二年相关条；军机处档、战事方略及地方奏报。编年记录适于钉住事件在朝廷文书中的时间位置，却需以地方、对方或物质材料检验其范围和后果。账本所列条件为：广西战场难以封堵太平军的机动路线。其后果被限定为：地方武装的出现改变后续战争权力结构。取证保留：以咸丰二年（1852年）的同期文书为定位起点；官修记录、奏报或外交文本服务特定机构立场，具体日期、规模、地方执行与对手视角须以独立材料复核。
\hypertarget{ledger:EQ-1852-HUNAN-ARMY-ORIGINS}{}\paragraph{咸丰二年（1852年）}曾国藩等在湖南组织团练和湘军雏形。核心取证为《清文宗实录》咸丰二年相关条；军机处档、战事方略及地方奏报。编年记录适于钉住事件在朝廷文书中的时间位置，却需以地方、对方或物质材料检验其范围和后果。账本所列条件为：八旗、绿营难以独力应对内战。其后果被限定为：地方募兵和筹饷为清廷提供新力量，也提高地方军政自主性。取证保留：以咸丰二年（1852年）的同期文书为定位起点；官修记录、奏报或外交文本服务特定机构立场，具体日期、规模、地方执行与对手视角须以独立材料复核。
\hypertarget{ledger:EQ-1853-NANJING-TAKING}{}\paragraph{咸丰三年（1853年）}太平军攻占南京并改名天京，建立长期政权中心。核心取证为《清文宗实录》咸丰三年相关条；军机处档、战事方略及地方奏报。编年记录适于钉住事件在朝廷文书中的时间位置，却需以地方、对方或物质材料检验其范围和后果。账本所列条件为：长江交通和清军防线崩解提供战略机会。其后果被限定为：江南财政文化中心易手，战争波及全国资源网络。取证保留：以咸丰三年（1853年）的同期文书为定位起点；官修记录、奏报或外交文本服务特定机构立场，具体日期、规模、地方执行与对手视角须以独立材料复核。
\hypertarget{ledger:EQ-1853-NORTHERN-EXPEDITION}{}\paragraph{咸丰三年（1853年）}太平军发动北伐，进逼华北但补给与指挥困难加重。核心取证为《清文宗实录》咸丰三年相关条；军机处档、战事方略及地方奏报。编年记录适于钉住事件在朝廷文书中的时间位置，却需以地方、对方或物质材料检验其范围和后果。账本所列条件为：天京建立后希望扩大战略压力。其后果被限定为：北伐失败限制太平天国对北方的直接威胁。取证保留：以咸丰三年（1853年）的同期文书为定位起点；官修记录、奏报或外交文本服务特定机构立场，具体日期、规模、地方执行与对手视角须以独立材料复核。
\hypertarget{ledger:EQ-1853-WESTERN-EXPEDITION}{}\paragraph{咸丰三年（1853年）}太平军西征争夺长江中上游，战事牵动安徽、江西、湖北。核心取证为《清文宗实录》咸丰三年相关条；军机处档、战事方略及地方奏报。编年记录适于钉住事件在朝廷文书中的时间位置，却需以地方、对方或物质材料检验其范围和后果。账本所列条件为：天京需要控制粮食、兵员和长江交通。其后果被限定为：区域战场长期化，地方社会承受征发与迁徙。取证保留：以咸丰三年（1853年）的同期文书为定位起点；官修记录、奏报或外交文本服务特定机构立场，具体日期、规模、地方执行与对手视角须以独立材料复核。
\hypertarget{ledger:EQ-1853-SMALL-SWORDS-SHANGHAI}{}\paragraph{咸丰三年（1853年）}上海小刀会起事，租界和华人城市空间关系发生变化。核心取证为《清文宗实录》咸丰三年相关条；地方志、督抚奏折及地方档案。编年记录适于钉住事件在朝廷文书中的时间位置，却需以地方、对方或物质材料检验其范围和后果。账本所列条件为：城市商业、秘密结社和地方治理矛盾交织。其后果被限定为：租界保护与地方镇压的边界需结合中外档案分析。取证保留：以咸丰三年（1853年）的同期文书为定位起点；官修记录、奏报或外交文本服务特定机构立场，具体日期、规模、地方执行与对手视角须以独立材料复核。
\hypertarget{ledger:EQ-1854-TAIPING-NORTHERN-DEFEAT}{}\paragraph{咸丰四年（1854年）}太平军北伐受挫，清军逐步恢复华北防线。核心取证为《清文宗实录》咸丰四年相关条；军机处档、战事方略及地方奏报。编年记录适于钉住事件在朝廷文书中的时间位置，却需以地方、对方或物质材料检验其范围和后果。账本所列条件为：长距离行军、冬季补给与地方支持不足限制北伐。其后果被限定为：太平天国转而依赖长江流域与地方战场。取证保留：以咸丰四年（1854年）的同期文书为定位起点；官修记录、奏报或外交文本服务特定机构立场，具体日期、规模、地方执行与对手视角须以独立材料复核。
\hypertarget{ledger:EQ-1854-XIANG-ARMY-ORGANIZATION}{}\paragraph{咸丰四年（1854年）}湘军编制、饷源和将领网络逐步成形。核心取证为《清文宗实录》咸丰四年相关条；宫中朱批奏折、内阁大库题本与《清史稿》相关志传。这类上谕、题本或部议首先证明机构处理了何种命令、呈报与责任，而不自动证明地方收到、执行或接受。账本所列条件为：清廷需要可用的地方军事力量。其后果被限定为：军事私人网络与国家授权之间的界线长期模糊。取证保留：以咸丰四年（1854年）的同期文书为定位起点；官修记录、奏报或外交文本服务特定机构立场，具体日期、规模、地方执行与对手视角须以独立材料复核。
\hypertarget{ledger:EQ-1855-YELLOW-RIVER-COURSE-SHIFT}{}\paragraph{咸丰五年（1855年）}黄河在铜瓦厢决口改道北流，重塑华北水系与社会经济。核心取证为《清文宗实录》咸丰五年相关条；地方志、督抚奏折及地方档案。编年记录适于钉住事件在朝廷文书中的时间位置，却需以地方、对方或物质材料检验其范围和后果。账本所列条件为：长期河工压力和水文变化共同作用。其后果被限定为：河道变迁的时间、影响范围和人口后果须依地方资料分区考察。取证保留：以咸丰五年（1855年）的同期文书为定位起点；官修记录、奏报或外交文本服务特定机构立场，具体日期、规模、地方执行与对手视角须以独立材料复核。
\hypertarget{ledger:EQ-1855-NIAN-REBELLION-EXPANSION}{}\paragraph{咸丰五年（1855年）}捻军在华北、安徽等地活动扩大，形成与太平战争交织的战场。核心取证为《清文宗实录》咸丰五年相关条；军机处档、战事方略及地方奏报。编年记录适于钉住事件在朝廷文书中的时间位置，却需以地方、对方或物质材料检验其范围和后果。账本所列条件为：灾荒、流民、盐枭和地方武装网络提供条件。其后果被限定为：捻军组织结构和社会基础不应由“匪”一词简化。取证保留：以咸丰五年（1855年）的同期文书为定位起点；官修记录、奏报或外交文本服务特定机构立场，具体日期、规模、地方执行与对手视角须以独立材料复核。
\hypertarget{ledger:EQ-1856-TIANJING-INCIDENT}{}\paragraph{咸丰六年（1856年）}天京事变中太平天国领导层发生内讧和大规模杀戮。核心取证为《清文宗实录》咸丰六年相关条；宫中朱批奏折、内阁大库题本与《清史稿》相关志传。这类上谕、题本或部议首先证明机构处理了何种命令、呈报与责任，而不自动证明地方收到、执行或接受。账本所列条件为：权力集中、军政分工与宗教权威冲突累积。其后果被限定为：事件严重削弱太平天国协调能力，细节和死亡数存在争议。取证保留：以咸丰六年（1856年）的同期文书为定位起点；官修记录、奏报或外交文本服务特定机构立场，具体日期、规模、地方执行与对手视角须以独立材料复核。
\hypertarget{ledger:EQ-1856-SECOND-OPIUM-WAR-BEGIN}{}\paragraph{咸丰六年（1856年）}亚罗号事件等冲突后英法对清开战，第二次鸦片战争开始。核心取证为《清文宗实录》咸丰六年相关条；军机处档、战事方略及地方奏报。编年记录适于钉住事件在朝廷文书中的时间位置，却需以地方、对方或物质材料检验其范围和后果。账本所列条件为：条约修订、领事权与商业扩张要求升级。其后果被限定为：战争与内战并行，迫使清廷分散军事和财政资源。取证保留：以咸丰六年（1856年）的同期文书为定位起点；官修记录、奏报或外交文本服务特定机构立场，具体日期、规模、地方执行与对手视角须以独立材料复核。
\hypertarget{ledger:EQ-1857-GUANGZHOU-FALL}{}\paragraph{咸丰七年（1857年）}英法联军攻占广州，叶名琛被俘。核心取证为《清文宗实录》咸丰七年相关条；军机处档、战事方略及地方奏报。编年记录适于钉住事件在朝廷文书中的时间位置，却需以地方、对方或物质材料检验其范围和后果。账本所列条件为：珠江口海军优势和清军防御薄弱。其后果被限定为：广东地方治理、商贸和对外关系受到持续冲击。取证保留：以咸丰七年（1857年）的同期文书为定位起点；官修记录、奏报或外交文本服务特定机构立场，具体日期、规模、地方执行与对手视角须以独立材料复核。
\hypertarget{ledger:EQ-1858-TIANJIN-TREATIES}{}\paragraph{咸丰八年（1858年）}《天津条约》等在战争压力下签订，扩大通商、使节与传教权利。核心取证为《清文宗实录》咸丰八年相关条；《筹办夷务始末》相关年卷与中外照会、条约文本。这类上谕、题本或部议首先证明机构处理了何种命令、呈报与责任，而不自动证明地方收到、执行或接受。账本所列条件为：联军北上威胁京师迫使清廷议约。其后果被限定为：签署不等于立即批准或实施，后续冲突仍在发生。取证保留：以咸丰八年（1858年）的同期文书为定位起点；官修记录、奏报或外交文本服务特定机构立场，具体日期、规模、地方执行与对手视角须以独立材料复核。
\hypertarget{ledger:EQ-1859-TAKU-BATTLE}{}\paragraph{咸丰九年（1859年）}大沽口战事使英法换约受阻，战争再度升级。核心取证为《清文宗实录》咸丰九年相关条；军机处档、战事方略及地方奏报。编年记录适于钉住事件在朝廷文书中的时间位置，却需以地方、对方或物质材料检验其范围和后果。账本所列条件为：条约批准和使节入京问题未获解决。其后果被限定为：清方短期胜利未改变总体军力与外交劣势。取证保留：以咸丰九年（1859年）的同期文书为定位起点；官修记录、奏报或外交文本服务特定机构立场，具体日期、规模、地方执行与对手视角须以独立材料复核。
\subsection{1860—1869：事件、文书与地方后果}
这一组节点按账本的日期顺序排列。相邻事件并不必然构成单一因果，却共同显示信息、资源与地方行动如何在同一时段交错。
\hypertarget{ledger:EQ-1860-BEIJING-CAMPAIGN}{}\paragraph{咸丰十年（1860年）}英法联军进逼北京，圆明园遭焚掠，咸丰帝出走热河。核心取证为《清文宗实录》咸丰十年相关条；军机处档、战事方略及地方奏报。编年记录适于钉住事件在朝廷文书中的时间位置，却需以地方、对方或物质材料检验其范围和后果。账本所列条件为：大沽后战争升级与清军防线崩溃。其后果被限定为：战争记忆、掠夺和城市影响需结合中外见证与物质证据。取证保留：以咸丰十年（1860年）的同期文书为定位起点；官修记录、奏报或外交文本服务特定机构立场，具体日期、规模、地方执行与对手视角须以独立材料复核。
\hypertarget{ledger:EQ-1860-CONVENTION-OF-PEKING}{}\paragraph{咸丰十年（1860年）}《北京条约》及相关中俄条约确认和扩大列强权益并涉及东北边界。核心取证为《清文宗实录》咸丰十年相关条；《筹办夷务始末》相关年卷与中外照会、条约文本。这类上谕、题本或部议首先证明机构处理了何种命令、呈报与责任，而不自动证明地方收到、执行或接受。账本所列条件为：北京陷落迫使清廷接受新安排。其后果被限定为：各条约文本、边界实施和土地控制不可混为一谈。取证保留：以咸丰十年（1860年）的同期文书为定位起点；官修记录、奏报或外交文本服务特定机构立场，具体日期、规模、地方执行与对手视角须以独立材料复核。
\hypertarget{ledger:EQ-1861-XIANFENG-DEATH}{}\paragraph{咸丰十一年（1861年）}咸丰帝去世，同治帝即位，顾命大臣与两宫太后权力冲突爆发。核心取证为《清文宗实录》咸丰十一年相关条；宫中朱批奏折、内阁大库题本与《清史稿》相关志传。这类上谕、题本或部议首先证明机构处理了何种命令、呈报与责任，而不自动证明地方收到、执行或接受。账本所列条件为：战时政局与幼主继承形成权力真空。其后果被限定为：辛酉政变重塑清廷中枢决策结构。取证保留：以咸丰十一年（1861年）的同期文书为定位起点；官修记录、奏报或外交文本服务特定机构立场，具体日期、规模、地方执行与对手视角须以独立材料复核。
\hypertarget{ledger:EQ-1861-XINYOU-COUP}{}\paragraph{咸丰十一年（1861年）}慈禧、慈安与恭亲王发动辛酉政变，肃顺等顾命大臣失势。核心取证为《清文宗实录》咸丰十一年相关条；宫中朱批奏折、内阁大库题本与《清史稿》相关志传。这类上谕、题本或部议首先证明机构处理了何种命令、呈报与责任，而不自动证明地方收到、执行或接受。账本所列条件为：幼主摄政安排与京师政治网络冲突。其后果被限定为：两宫垂帘与恭亲王参与中枢成为同治初年格局。取证保留：以咸丰十一年（1861年）的同期文书为定位起点；官修记录、奏报或外交文本服务特定机构立场，具体日期、规模、地方执行与对手视角须以独立材料复核。
\hypertarget{ledger:EQ-1861-ZONGLI-YAMEN}{}\paragraph{咸丰十一年（1861年）}总理各国事务衙门设立，清廷开始形成常设近代外交处理机构。核心取证为《清文宗实录》咸丰十一年相关条；《筹办夷务始末》相关年卷与中外照会、条约文本。这类上谕、题本或部议首先证明机构处理了何种命令、呈报与责任，而不自动证明地方收到、执行或接受。账本所列条件为：条约体系和驻京使节要求持续对外交涉能力。其后果被限定为：机构职责不断调整，不能视为完整现代外交部。取证保留：以咸丰十一年（1861年）的同期文书为定位起点；官修记录、奏报或外交文本服务特定机构立场，具体日期、规模、地方执行与对手视角须以独立材料复核。
\hypertarget{ledger:EQ-1862-SELF-STRENGTHENING-ORIGINS}{}\paragraph{同治元年（1862年）}洋务实践在军械、翻译、外交和地方军政中展开。核心取证为《清穆宗实录》同治元年相关条；宫中朱批奏折、内阁大库题本与《清史稿》相关志传。这类上谕、题本或部议首先证明机构处理了何种命令、呈报与责任，而不自动证明地方收到、执行或接受。账本所列条件为：内战与列强压力促使官员寻求技术和制度调整。其后果被限定为：不同地区的洋务项目具有不同资金、目标和绩效。取证保留：以同治元年（1862年）的同期文书为定位起点；官修记录、奏报或外交文本服务特定机构立场，具体日期、规模、地方执行与对手视角须以独立材料复核。
\hypertarget{ledger:EQ-1862-SHANGHAI-FOREIGN-ARMS}{}\paragraph{同治元年（1862年）}上海的军火、雇佣力量和租界网络影响清军对太平军作战。核心取证为《清穆宗实录》同治元年相关条；军机处档、战事方略及地方奏报。编年记录适于钉住事件在朝廷文书中的时间位置，却需以地方、对方或物质材料检验其范围和后果。账本所列条件为：江南战场接近条约口岸与国际商业网络。其后果被限定为：外来军事技术和人员角色须避免夸大为决定性单因。取证保留：以同治元年（1862年）的同期文书为定位起点；官修记录、奏报或外交文本服务特定机构立场，具体日期、规模、地方执行与对手视角须以独立材料复核。
\hypertarget{ledger:EQ-1863-EVER-VICTORIOUS-ARMY}{}\paragraph{同治二年（1863年）}常胜军在江南战场参与清军作战，戈登接掌后扩大影响。核心取证为《清穆宗实录》同治二年相关条；军机处档、战事方略及地方奏报。编年记录适于钉住事件在朝廷文书中的时间位置，却需以地方、对方或物质材料检验其范围和后果。账本所列条件为：上海周边的财政、外商和地方军政资源可被军事化。其后果被限定为：常胜军规模有限，其作用需放入湘淮军和地方战场比较。取证保留：以同治二年（1863年）的同期文书为定位起点；官修记录、奏报或外交文本服务特定机构立场，具体日期、规模、地方执行与对手视角须以独立材料复核。
\hypertarget{ledger:EQ-1863-NIAN-WAR-MOBILIZATION}{}\paragraph{同治二年（1863年）}清廷调集湘淮军和地方力量应对捻军，华北战事加剧。核心取证为《清穆宗实录》同治二年相关条；军机处档、战事方略及地方奏报。编年记录适于钉住事件在朝廷文书中的时间位置，却需以地方、对方或物质材料检验其范围和后果。账本所列条件为：太平战场与北方流动战场同时消耗军费。其后果被限定为：军队调动改变省际财政与地方权力关系。取证保留：以同治二年（1863年）的同期文书为定位起点；官修记录、奏报或外交文本服务特定机构立场，具体日期、规模、地方执行与对手视角须以独立材料复核。
\hypertarget{ledger:EQ-1864-TIANJING-FALL}{}\paragraph{同治三年（1864年）}湘军攻陷天京，太平天国核心政权覆亡。核心取证为《清穆宗实录》同治三年相关条；军机处档、战事方略及地方奏报。编年记录适于钉住事件在朝廷文书中的时间位置，却需以地方、对方或物质材料检验其范围和后果。账本所列条件为：长期围困、粮食危机和内部裂解削弱天京。其后果被限定为：战后屠杀、人口损失和城市重建须以多种材料谨慎处理。取证保留：以同治三年（1864年）的同期文书为定位起点；官修记录、奏报或外交文本服务特定机构立场，具体日期、规模、地方执行与对手视角须以独立材料复核。
\hypertarget{ledger:EQ-1865-DUNGAN-REVOLT}{}\paragraph{同治四年（1865年）}西北回民起事在陕西、甘肃等地扩展，形成长期战争。核心取证为《清穆宗实录》同治四年相关条；军机处档、战事方略及地方奏报。编年记录适于钉住事件在朝廷文书中的时间位置，却需以地方、对方或物质材料检验其范围和后果。账本所列条件为：地方冲突、军政失序和族群分类相互激化。其后果被限定为：动机、伤亡和社区边界须用地方与多语材料审慎叙述。取证保留：以同治四年（1865年）的同期文书为定位起点；官修记录、奏报或外交文本服务特定机构立场，具体日期、规模、地方执行与对手视角须以独立材料复核。
\hypertarget{ledger:EQ-1865-NIAN-LEADER-DEATH}{}\paragraph{同治四年（1865年）}张乐行等捻军领导人先后失势或被杀，捻军仍持续活动。核心取证为《清穆宗实录》同治四年相关条；军机处档、战事方略及地方奏报。编年记录适于钉住事件在朝廷文书中的时间位置，却需以地方、对方或物质材料检验其范围和后果。账本所列条件为：清军加强围堵与地方武装协作。其后果被限定为：领导层变化不等于流动作战网络立即瓦解。取证保留：以同治四年（1865年）的同期文书为定位起点；官修记录、奏报或外交文本服务特定机构立场，具体日期、规模、地方执行与对手视角须以独立材料复核。
\hypertarget{ledger:EQ-1866-FUJIAN-SHIPYARD}{}\paragraph{同治五年（1866年）}福州船政局创办，成为洋务造船、教育与技术翻译的重要项目。核心取证为《清穆宗实录》同治五年相关条；户部奏销、盐政、河工或海关档案。这类上谕、题本或部议首先证明机构处理了何种命令、呈报与责任，而不自动证明地方收到、执行或接受。账本所列条件为：海防压力和技术学习需求增长。其后果被限定为：厂局经费、产能和军事效果需用账册与具体船只记录评估。取证保留：以同治五年（1866年）的同期文书为定位起点；官修记录、奏报或外交文本服务特定机构立场，具体日期、规模、地方执行与对手视角须以独立材料复核。
\hypertarget{ledger:EQ-1866-NINGBO-SHANGHAI-INDUSTRY}{}\paragraph{同治五年（1866年）}江南制造局等军工项目在战后江南环境中继续发展。核心取证为《清穆宗实录》同治五年相关条；户部奏销、盐政、河工或海关档案。这类上谕、题本或部议首先证明机构处理了何种命令、呈报与责任，而不自动证明地方收到、执行或接受。账本所列条件为：湘淮军和条约口岸带来技术、资金与人才条件。其后果被限定为：“自强”口号不能替代对生产规模和依赖进口的检验。取证保留：以同治五年（1866年）的同期文书为定位起点；官修记录、奏报或外交文本服务特定机构立场，具体日期、规模、地方执行与对手视角须以独立材料复核。
\hypertarget{ledger:EQ-1867-NIAN-CAVALRY-CAMPAIGN}{}\paragraph{同治六年（1867年）}清军以河防、堡垒和机动军队围堵捻军骑兵。核心取证为《清穆宗实录》同治六年相关条；军机处档、战事方略及地方奏报。编年记录适于钉住事件在朝廷文书中的时间位置，却需以地方、对方或物质材料检验其范围和后果。账本所列条件为：捻军流动作战突破传统省界防御。其后果被限定为：军事策略的效果应与地方破坏和民众负担并看。取证保留：以同治六年（1867年）的同期文书为定位起点；官修记录、奏报或外交文本服务特定机构立场，具体日期、规模、地方执行与对手视角须以独立材料复核。
\hypertarget{ledger:EQ-1868-NIAN-SUPPRESSION}{}\paragraph{同治七年（1868年）}东捻军等主力被击败，捻军大规模战争告终。核心取证为《清穆宗实录》同治七年相关条；军机处档、战事方略及地方奏报。编年记录适于钉住事件在朝廷文书中的时间位置，却需以地方、对方或物质材料检验其范围和后果。账本所列条件为：淮军等部队的围堵和战场分割发挥作用。其后果被限定为：战后地方武装和财政动员遗产继续存在。取证保留：以同治七年（1868年）的同期文书为定位起点；官修记录、奏报或外交文本服务特定机构立场，具体日期、规模、地方执行与对手视角须以独立材料复核。
\hypertarget{ledger:EQ-1869-YANGZHOU-SALT-CRISIS}{}\paragraph{同治八年（1869年）}两淮盐政改革、商欠与财政重整成为中兴时期的重要议题。核心取证为《清穆宗实录》同治八年相关条；户部奏销、盐政、河工或海关档案。这类上谕、题本或部议首先证明机构处理了何种命令、呈报与责任，而不自动证明地方收到、执行或接受。账本所列条件为：战乱破坏盐业网络，厘金与新税制改变利益结构。其后果被限定为：改革效果需比较名义制度、实收与商人账目。取证保留：以同治八年（1869年）的同期文书为定位起点；官修记录、奏报或外交文本服务特定机构立场，具体日期、规模、地方执行与对手视角须以独立材料复核。
\subsection{1870—1879：事件、文书与地方后果}
这一组节点按账本的日期顺序排列。相邻事件并不必然构成单一因果，却共同显示信息、资源与地方行动如何在同一时段交错。
\hypertarget{ledger:EQ-1870-TIANJIN-MASSACRE}{}\paragraph{同治九年（1870年）}天津教案发生，法国领事、修女和中国民众死伤，引发外交危机。核心取证为《清穆宗实录》同治九年相关条；《筹办夷务始末》相关年卷与中外照会、条约文本。这类上谕、题本或部议首先证明机构处理了何种命令、呈报与责任，而不自动证明地方收到、执行或接受。账本所列条件为：教会活动、地方谣言和条约保护关系高度紧张。其后果被限定为：责任叙述在中法档案中差异显著，须避免单一归因。取证保留：以同治九年（1870年）的同期文书为定位起点；官修记录、奏报或外交文本服务特定机构立场，具体日期、规模、地方执行与对手视角须以独立材料复核。
\hypertarget{ledger:EQ-1870-ZHILI-FAMINE-RELIEF}{}\paragraph{同治九年（1870年）}华北灾荒和赈务问题促使官员、士绅和国际救济网络介入。核心取证为《清穆宗实录》同治九年相关条；地方志、督抚奏折及地方档案。编年记录适于钉住事件在朝廷文书中的时间位置，却需以地方、对方或物质材料检验其范围和后果。账本所列条件为：旱灾、粮价、交通和地方储备共同影响灾情。其后果被限定为：慈善报告与官府奏报的覆盖范围和数字需核对。取证保留：以同治九年（1870年）的同期文书为定位起点；官修记录、奏报或外交文本服务特定机构立场，具体日期、规模、地方执行与对手视角须以独立材料复核。
\hypertarget{ledger:EQ-1871-MUDAN-INCIDENT}{}\paragraph{同治十年（1871年）}琉球船民在台湾牡丹社一带遇害，事件成为日本对台政策借口。核心取证为《清穆宗实录》同治十年相关条；《筹办夷务始末》相关年卷与中外照会、条约文本。这类上谕、题本或部议首先证明机构处理了何种命令、呈报与责任，而不自动证明地方收到、执行或接受。账本所列条件为：海难、原住民地域与清朝地方治理边界交错。其后果被限定为：事件叙述须区分原住民行动、清方管辖与日方外交利用。取证保留：以同治十年（1871年）的同期文书为定位起点；官修记录、奏报或外交文本服务特定机构立场，具体日期、规模、地方执行与对手视角须以独立材料复核。
\hypertarget{ledger:EQ-1871-QING-STUDENTS-ABROAD}{}\paragraph{同治十年（1871年）}清廷派遣幼童赴美留学的计划启动，洋务教育进入海外培训阶段。核心取证为《清穆宗实录》同治十年相关条；内阁档、文集或报刊相关记录。编年记录适于钉住事件在朝廷文书中的时间位置，却需以地方、对方或物质材料检验其范围和后果。账本所列条件为：技术、语言与外交人才需求上升。其后果被限定为：选派和留学经历不代表全国教育制度已转型。取证保留：以同治十年（1871年）的同期文书为定位起点；官修记录、奏报或外交文本服务特定机构立场，具体日期、规模、地方执行与对手视角须以独立材料复核。
\hypertarget{ledger:EQ-1872-CHINA-MERCHANTS-STEAM-NAVIGATION}{}\paragraph{同治十一年（1872年）}轮船招商局创办，试图以官督商办方式发展近代航运。核心取证为《清穆宗实录》同治十一年相关条；户部奏销、盐政、河工或海关档案。这类上谕、题本或部议首先证明机构处理了何种命令、呈报与责任，而不自动证明地方收到、执行或接受。账本所列条件为：外轮竞争和沿海、长江运输需求推动新企业。其后果被限定为：企业性质、资本来源和经营绩效需依账册分析。取证保留：以同治十一年（1872年）的同期文书为定位起点；官修记录、奏报或外交文本服务特定机构立场，具体日期、规模、地方执行与对手视角须以独立材料复核。
\hypertarget{ledger:EQ-1873-TONGZHI-EMPEROR-REIGN}{}\paragraph{同治十二年（1873年）}同治帝亲政，垂帘体制与皇帝个人政治角色发生调整。核心取证为《清穆宗实录》同治十二年相关条；宫中朱批奏折、内阁大库题本与《清史稿》相关志传。这类上谕、题本或部议首先证明机构处理了何种命令、呈报与责任，而不自动证明地方收到、执行或接受。账本所列条件为：幼帝成年改变摄政安排。其后果被限定为：实际决策仍受太后、军机和重臣网络制约。取证保留：以同治十二年（1873年）的同期文书为定位起点；官修记录、奏报或外交文本服务特定机构立场，具体日期、规模、地方执行与对手视角须以独立材料复核。
\hypertarget{ledger:EQ-1873-JAPAN-RYUKYU-DISPUTE}{}\paragraph{同治十二年（1873年）}日本围绕琉球与台湾事务加强对清交涉。核心取证为《清穆宗实录》同治十二年相关条；《筹办夷务始末》相关年卷与中外照会、条约文本。这类上谕、题本或部议首先证明机构处理了何种命令、呈报与责任，而不自动证明地方收到、执行或接受。账本所列条件为：日本明治国家扩张与东亚朝贡关系重组。其后果被限定为：册封、属国和主权词汇应置于各方文本语境中。取证保留：以同治十二年（1873年）的同期文书为定位起点；官修记录、奏报或外交文本服务特定机构立场，具体日期、规模、地方执行与对手视角须以独立材料复核。
\hypertarget{ledger:EQ-1874-JAPANESE-TAIWAN-EXPEDITION}{}\paragraph{同治十三年（1874年）}日本以牡丹社事件为由出兵台湾南部，清廷被迫应对。核心取证为《清穆宗实录》同治十三年相关条；军机处档、战事方略及地方奏报。编年记录适于钉住事件在朝廷文书中的时间位置，却需以地方、对方或物质材料检验其范围和后果。账本所列条件为：清廷对台湾东部和原住民地区的治理能力有限。其后果被限定为：战后赔款与外交安排刺激清廷加强台湾经营。取证保留：以同治十三年（1874年）的同期文书为定位起点；官修记录、奏报或外交文本服务特定机构立场，具体日期、规模、地方执行与对手视角须以独立材料复核。
\hypertarget{ledger:EQ-1874-TONGZHI-DEATH}{}\paragraph{同治十三年（1874年）}同治帝去世，慈禧选择光绪帝继位，再次形成幼主政治。核心取证为《清穆宗实录》同治十三年相关条；宫中朱批奏折、内阁大库题本与《清史稿》相关志传。这类上谕、题本或部议首先证明机构处理了何种命令、呈报与责任，而不自动证明地方收到、执行或接受。账本所列条件为：同治无嗣使宗法与宫廷权力问题交织。其后果被限定为：慈禧重新垂帘，晚清中枢进入新阶段。取证保留：以同治十三年（1874年）的同期文书为定位起点；官修记录、奏报或外交文本服务特定机构立场，具体日期、规模、地方执行与对手视角须以独立材料复核。
